%=============================================================================
% 前人主要研究成果综述
%=============================================================================

\section*{三、国内外主要研究成果}

\subsection*{1. 反应堆中微子拟合程序的高性能拟合算法研究}

在反应堆中微子振荡的精确测量中，Feldman--Cousins 覆盖率检验等统计方法要求对百万量级的伪实验进行完整拟合，其计算瓶颈主要集中在能谱的前向折叠（forward-folding）构造过程——包括拟合过程中探测器响应矩阵的反复求值与 IBD 微分截面映射的大规模矩阵运算。针对这一挑战，前人已发展了若干高效优化算法。GooStats 框架~\cite{Ding_2018} 提出了预载探测器响应网格（Response Grid）的策略，通过将耗时的响应计算在拟合前一次性完成并常驻内存，避免了每次迭代中的重复求值，实现了超过一个数量级的运算加速。ReMU 框架~\cite{Koch2019} 则从矩阵存储角度出发，采用稀疏矩阵表示，仅保留含非零元素的行，并结合高效线性代数操作显著提升了能谱构造效率。本课题在上述"预计算"与"稀疏化"思想的基础上，进一步发展了基于张量计算（Tensor-based）的高性能拟合框架，将反应堆能谱生成、IBD 映射与探测器响应统一构建为张量运算链，结合结果缓存、带状稀疏矩阵与分块卷积三项核心优化，在保持数值精度的前提下实现了约一个数量级的综合加速。

\subsection*{2. 中微子实验的远近点联合分析与能谱约束策略}

利用近点探测器数据约束谱型系统误差是中微子振荡实验的核心分析策略之一。在加速器中微子实验中，T2K~\cite{Abe2020T2K} 与 NOvA~\cite{Acero2022ProfiledFC} 通过远近点探测器同步测量，利用近点数据对通量归一化与相互作用截面模型施加强约束，从而有效抵消理论模型依赖，使振荡参数的提取更为稳健。在 JUNO 早期运行阶段，借鉴了上述加速器实验的联合分析逻辑，在近探测器 TAO 尚未积累充足统计量的条件下，将大亚湾（Daya Bay）实验已积累的超过 550 万个反中微子事例的高统计量近点能谱数据引入 JUNO 远场振荡分析，作为反应堆能谱形状的外部约束。针对大亚湾能量分辨率低以及 JUNO 与大亚湾的探测器响应存在差异，如何在统计误差与模型依赖的系统误差之间寻找平衡是核心挑战。本课题正是在上述的联合分析框架之上，侧重于分bin策略的优化，通过大量 Toy-MC 模拟，比较了固定反应堆模型与包含精细结构扰动的变化模型两种情形，在不同段宽度下振荡参数精度与偏差的表现，量化反应堆模型依赖引入的系统误差，确定最优的分段策略。


\subsection*{3. 复杂参数空间下的 Feldman-Cousins 方法应用}

在中微子振荡参数的置信区间构造中，Wilks 定理~\cite{Wilks:1938dza} 所依赖的渐近正态假设在低统计量、参数接近物理边界或似然函数高度非线性等情形下可能失效，导致基于固定 $\Delta\chi^2$ 阈值的置信区间出现覆盖率不足。为解决这一问题，Feldman 与 Cousins~\cite{Feldman_1998} 提出了基于似然比排序与 Neyman 置信带构造的频率学方法，通过大量伪实验直接构建检验统计量的经验分布，从而无需依赖渐近近似即可保证正确的覆盖性质。该方法已在多个中微子实验中得到成功应用。NOvA 实验~\cite{Acero2022ProfiledFC} 在振荡参数测量中发现，当参数接近物理边界或存在多重局部极小值时，$\Delta\chi^2$ 分布显著偏离 Wilks 预期的卡方分布，因此系统性地采用了 Profiled Feldman-Cousins 方法以保证严格的频率学覆盖率。本课题也采用Feldman-Cousins方法，针对 JUNO 低统计量条件下 $\Delta m^2_{31}$ 似然面的多峰结构与非抛物线行为，运用 Feldman-Cousins 方法进行覆盖率检验，确保置信区间的统计稳健性。
